\chapter{实验设计与结果分析}


\section{实验数据与实验设置}

\subsection{数据集与预处理}

本文采用自然场景数据集（Natural Scenes Dataset, NSD）进行模型训练与评估。NSD 是当前规模较大且标注完善的视觉神经数据集之一，包含来自 8 名被试在观看自然图像时采集的功能磁共振成像（fMRI）信号。每位被试共观看约 8{,}000 至 9{,}000 张图像，所有刺激图像均来源于 MS-COCO 数据集，涵盖丰富的现实世界场景与多样化语义内容。

在数据使用上，本文直接采用 NSD 官方发布的预处理后的皮层顶点级（vertex-wise）fMRI 响应数据。该数据已完成标准化预处理流程，包括空间对齐、去噪以及信号标准化等，从而保证不同被试之间数据具有良好的可比性。此外，所有视觉刺激均与对应的神经响应严格配对，用于监督模型学习跨模态映射关系。

在实验设置上，本文采用严格的零样本跨被试（zero-shot cross-subject）评估协议。具体而言，模型在多个源被试的数据上进行训练，并在完全未参与训练的目标被试上直接测试。在整个测试过程中，不进行任何形式的被试特定微调或校准，从而真实评估模型的跨被试泛化能力。实验结果在多个目标被试上进行统计，并报告平均性能。

\subsection{对比方法}

为了全面评估本文提出的 MINDMELD 方法的有效性，本文选取了当前具有代表性的脑信号视觉重建方法作为对比基线，涵盖完全监督、少样本以及零样本三种设置：

\begin{itemize}
    \item \textbf{MindEye2}：基于扩散模型的视觉重建方法，能够通过强大的生成能力实现高质量图像重建，在完全微调和少样本设置下表现优异。

    \item \textbf{NeuroPictor}：通过学习神经信号与语义空间之间的映射关系进行图像重建的方法，在语义表达方面具有较强能力。本文采用其零样本版本（NeuroPictor*），即仅在源被试上训练后直接测试。

    \item \textbf{ZEBRA}：专注于跨被试对齐问题的方法，通过构建跨被试一致的表示空间来提升泛化能力，同时包含一个不带对齐策略的 baseline 版本。

    \item \textbf{MindTuner}：通过少量目标被试数据进行快速适配的方法，在 few-shot 设置下具有较强表现。

    \item 其他代表性方法（如 Takagi 等、Ozcelik 等、MindBridge、UMBRAE 等）主要用于提供完全监督条件下的性能上界参考。
\end{itemize}

需要说明的是，为保证对比公平性，所有零样本方法均严格遵循相同的训练与测试划分，不使用任何目标被试信息。

\subsection{评价指标}

为了从不同层面全面评估模型的重建性能，本文采用低层视觉指标与高层语义指标两类评价标准。

\subsubsection{低层视觉指标}

低层指标主要用于衡量重建图像在像素空间和结构层面的相似性，包括：

\begin{itemize}
    \item \textbf{PixCorr}：像素级相关性指标，用于衡量重建图像与原始图像在像素层面的线性相关性。

    \item \textbf{SSIM}：结构相似性指标，从亮度、对比度和结构三个方面评估图像相似性。

    \item \textbf{Alex(2) / Alex(5)}：基于 AlexNet 不同层特征的相似度指标，用于衡量重建图像在中低层视觉特征空间中的一致性。
\end{itemize}

\subsubsection{高层语义指标}

高层指标用于评估重建图像在语义层面的表达能力和感知一致性，包括：

\begin{itemize}
    \item \textbf{Incep}：基于 Inception 网络特征的语义一致性指标。

    \item \textbf{CLIP}：利用 CLIP 模型评估重建图像与其语义表示之间的一致性。

    \item \textbf{Eff}：基于 EfficientNet 特征的感知距离指标（数值越低表示越接近）。

    \item \textbf{SwAV}：基于自监督学习特征的语义距离指标（数值越低表示越接近）。
\end{itemize}

所有方法均在统一的评价协议下进行评估，以确保不同方法之间的比较具有公平性和一致性。


\section{整体重建结果对比}

\begin{table*}[t]
  \centering
  \setlength{\tabcolsep}{3pt}
  \footnotesize
  \caption{MINDMELD 与代表性方法在 NSD 数据集上的整体重建结果对比。}
  \label{tab:main_results}
  \begin{tabular}{lcccccccc}
    \toprule
    方法 & PixCorr $\uparrow$ & SSIM $\uparrow$ & Alex(2) $\uparrow$ & Alex(5) $\uparrow$ & Incep $\uparrow$ & CLIP $\uparrow$ & Eff $\downarrow$ & SwAV $\downarrow$ \\
    \midrule
    \multicolumn{9}{c}{\textbf{全量微调}} \\
    \midrule
    Takagi et al. & 0.246 & 0.410 & 78.9\% & 85.6\% & 83.8\% & 82.1\% & 0.811 & 0.504 \\
    Ozcelik et al. & 0.273 & 0.365 & 94.4\% & 96.6\% & 91.3\% & 90.9\% & 0.728 & 0.422 \\
    MindEye1 & 0.319 & 0.360 & 92.8\% & 96.9\% & 94.6\% & 93.3\% & 0.648 & 0.377 \\
    UMBRAE & 0.283 & 0.341 & 95.5\% & 97.0\% & 91.7\% & 93.5\% & 0.700 & 0.393 \\
    NeuroPictor & 0.229 & 0.375 & 96.5\% & 98.4\% & 94.5\% & 93.3\% & 0.639 & 0.350 \\
    MindBridge & 0.151 & 0.263 & 87.7\% & 95.5\% & 92.4\% & 94.7\% & 0.712 & 0.418 \\
    MindEye2 & 0.322 & 0.431 & 96.1\% & 98.6\% & 95.4\% & 93.0\% & 0.619 & 0.344 \\
    MindTuner & 0.322 & 0.421 & 95.8\% & 98.8\% & 95.6\% & 93.8\% & 0.612 & 0.340 \\
    \midrule
    \multicolumn{9}{c}{\textbf{少样本迁移}} \\
    \midrule
    MindEye2（1小时） & 0.195 & 0.419 & 84.2\% & 90.6\% & 81.2\% & 79.2\% & 0.810 & 0.468 \\
    MindTuner（1小时） & 0.224 & 0.420 & 87.8\% & 93.6\% & 84.8\% & 83.5\% & 0.780 & 0.440 \\
    \midrule
    \multicolumn{9}{c}{\textbf{零样本迁移}} \\
    \midrule
    NeuroPictor* & 0.057 & 0.297 & 71.4\% & 74.7\% & 62.5\% & 66.0\% & 0.939 & 0.607 \\
    ZEBRA baseline & 0.074 & 0.316 & 70.8\% & 74.0\% & 63.5\% & 62.5\% & 0.920 & 0.602 \\
    ZEBRA & \textbf{0.131} & 0.375 & 74.6\% & \textbf{81.2\%} & 72.2\% & 71.5\% & 0.837 & 0.506 \\
    \textbf{MINDMELD（ours）} & \underline{0.119} & \textbf{0.405} & \textbf{75.0\%} & \underline{80.3\%} & \textbf{75.6\%} & \textbf{73.2\%} & \textbf{0.816} & \textbf{0.443} \\
    \bottomrule
  \end{tabular}
\end{table*}

\subsection{定量结果比较}

表~\ref{tab:main_results} 展示了在 NSD 数据集上的定量比较结果，涵盖完全微调、少样本以及严格零样本三种设置。所有结果在被试 1、2、5 和 7 上取平均，以保证评估的稳定性。

在零样本跨被试设置下，MINDMELD 在多个关键指标上取得了当前最优或具有竞争力的结果。在低层视觉指标方面，MINDMELD 在 SSIM 上取得最优结果（0.405），相比 ZEBRA（0.375）有明显提升，说明其在图像结构恢复方面具有更强能力。在 PixCorr 指标上，MINDMELD（0.119）略低于 ZEBRA（0.131），但仍显著优于其他零样本方法，体现出较强的像素级一致性。

在特征层指标方面，MINDMELD 在 Alex(2) 上取得最佳结果（75.0\%），在 Alex(5) 上也接近最优方法，说明其在中低层视觉特征空间中能够更好地对齐真实图像。

在高层语义指标方面，MINDMELD 表现出显著优势。在 Incep（75.6\%）和 CLIP（73.2\%）指标上均优于所有对比方法，表明模型在语义表达与跨模态对齐方面具有更强能力。同时，在 Eff（0.816）和 SwAV（0.443）两个距离型指标上，MINDMELD 也取得最优结果，说明其重建结果在高层特征空间中更加接近真实图像。

与少样本方法相比，尽管这些方法利用了目标被试的少量数据进行适配，MINDMELD 在多个指标上仍具有竞争力，说明其在无需目标被试参与的情况下具备良好的泛化能力。

\subsection{定性重建结果比较}

图~\ref{fig:sota} 展示了不同方法在零样本跨被试设置下的可视化对比结果。可以观察到，MINDMELD 能够较为准确地恢复图像的主体语义、场景结构以及关键视觉属性。在包含复杂背景或多物体的场景中，本文方法仍能保持较好的整体布局一致性。

相比之下，部分基线方法虽然能够恢复局部颜色或纹理信息，但在主体识别和语义一致性方面存在明显偏差，例如出现语义混淆或结构失真的情况。

此外，在细节表现方面，MINDMELD 在边界清晰度和结构完整性上表现更优，生成结果整体更加自然。这说明本文方法学习到的被试不变语义表示能够更稳定地指导生成模型进行图像重建。

\subsection{结果小结}

综合定量与定性结果可以得出，MINDMELD 在严格零样本跨被试设置下，在低层结构保真度与高层语义一致性之间实现了更优的平衡。尤其是在多个语义指标和感知距离指标上的提升，表明本文方法在跨被试语义建模方面具有明显优势。

总体而言，实验结果验证了本文所提出方法的有效性，即通过神经响应扩展与语义解耦机制，可以在不依赖目标被试数据的前提下显著提升视觉重建性能。


\section{消融实验与模块有效性分析}

\subsection{被试感知神经响应预测模块的作用}

\subsection{语义解耦结构的作用}

\subsection{对抗约束与一致性约束的作用}


\section{表示分析与可解释性}

\subsection{皮层层级结构的自发涌现}

模型学习到的顶点级层偏好能够自发恢复符合神经科学认知的皮层层级结构。在没有引入任何显式解剖先验、ROI 约束或层级监督信号的情况下，模型在每个皮层顶点上学习到对 RADIOv4 不同深度特征的结构化偏好，从而在整体上形成由浅至深的层级对齐模式。这一现象与人类视觉系统从低级视觉皮层（如 V1/V2）到高级语义区域逐级抽象的经典层级组织高度一致，说明模型在纯数据驱动的条件下，能够自动捕捉神经表征中的层级结构。

从机制上看，这种层级涌现反映了不同深度视觉特征在解释脑响应时的功能分工：浅层特征更偏向编码局部纹理、边缘与低级统计信息，而深层特征则更侧重语义与整体概念表达。模型通过对这些多尺度特征的自适应加权，实现了对神经响应的高效建模。

进一步分析发现一个稳定且具有启发性的现象：最浅层与最深层特征从未在任何顶点上占据主导地位，但在几乎整个皮层范围内始终保持非零且稳定的低权重。这表明，这些极端层级并非直接决定神经响应的主要来源，而更可能承担“全局约束”或“上下文补充”的作用。具体而言，浅层特征可能提供基础的感知约束以维持局部一致性，而深层特征则提供粗粒度语义背景，从而辅助中间层特征进行表达。相比之下，中间层特征在多数顶点上占据主导地位，说明其在连接低级视觉结构与高级语义之间起到了关键桥梁作用。

这一结果不仅验证了模型设计的合理性，也表明所提出的方法能够在无监督结构约束的情况下，学习到具有神经科学意义的可解释表示。

\subsection{语义表示的被试不变性与判别性}

学习到的语义 token 构建了一个被试不变且具有语义表达能力的潜在空间。为了验证模型是否成功实现“语义—被试”解耦，我们对语义表示 $\mathbf{Z}_{\mathrm{sem}}$ 的分布结构进行了系统分析。

从被试维度来看，不同个体的数据在该潜在空间中呈现高度混合状态，说明表示中已基本消除了与个体相关的特征信息。这一现象表明，对抗解耦机制有效抑制了被试特异性信号的残留，从而避免模型依赖个体差异进行预测。这对于跨被试任务尤为关键，因为任何残留的个体信息都可能导致泛化能力下降。

与此同时，从语义维度观察，该潜在空间呈现出清晰的类别结构，不同语义类别之间具有良好的分离性。这说明模型在去除被试信息的同时，并未损失关键语义信息，反而保留并强化了与视觉内容相关的判别特征。

这种“被试混合 + 语义分离”的结构具有重要意义：它表明模型成功将神经响应分解为“共享语义成分”与“个体差异成分”，并在语义空间中仅保留前者。这一性质正是实现零样本跨被试重建的核心前提，因为只有当表示空间对不同被试具有一致语义解释时，模型才能在未见过目标被试的情况下进行有效推断。

从表示学习的角度看，这种结构可以理解为一种受约束的判别性嵌入：既满足跨域不变性，又保留任务相关信息，从而在泛化能力与表达能力之间取得平衡。

\subsection{跨被试语义对齐机制分析}

跨重建目标有效实现了跨被试语义对齐。为了进一步理解模型如何实现跨被试泛化，我们对跨重建机制的作用进行了分析。结果表明，当不同被试观看相同刺激时，其对应的潜在表示会在训练过程中逐渐收敛到一个共享的语义中心。这说明跨重建损失在显式地约束表示空间，使其对相同语义输入产生一致响应。

与传统方法依赖隐式对齐不同，这种设计通过直接在表示层施加一致性约束，使模型能够主动消除个体差异带来的偏移，从而在语义层面实现统一建模。换言之，模型学习的并非“被试间映射关系”，而是一个对所有被试通用的语义坐标系。

这一机制的优势在于，其对齐目标具有明确的语义意义，因此比单纯的分布匹配或对抗训练更加稳定且可控。同时，由于对齐是在语义层完成的，模型能够在不依赖低层精细对齐的情况下，实现跨被试的信息整合。

值得注意的是，该对齐过程并未破坏表示的生成兼容性。尽管模型没有显式引入低层重建分支，但所学习到的语义表示仍能够与图像生成模块（如 IP-Adapter Plus）有效协同，在保持语义一致性的同时恢复合理的空间结构。这说明该表示在抽象语义与空间信息之间保持了良好的信息平衡。

\subsection{本节小结}

综合上述分析可以看出，MINDMELD 在表示层面具备良好的结构性与可解释性：首先，模型能够自发学习到符合生物视觉系统的层级组织；其次，其语义表示在实现被试不变性的同时，仍保持较强的判别能力；最后，通过跨重建机制，模型在语义空间中实现了稳定的跨被试对齐。

这些性质共同构成了模型实现零样本跨被试重建的关键基础，也从机制层面解释了其在定量与定性实验中取得性能提升的原因。


\chapterbib